\section{Control flow and operating states}
\begin{figure}[htbp]\centering\begin{tikzpicture}[>=Stealth,node distance=7mm,box/.style={draw=SEDSBlue,rounded corners,align=center,text width=110mm,minimum height=10mm,font=\small}]
\node[box] (hal) {Reset: HAL / clocks / GPIO selects high / SPI1 / USB PCD};
\node[box,below=of hal] (rtos) {Start ThreadX; create queues, mutex, events and static stacks};
\node[box,below=of rtos] (check) {Sensor identity/config checks; data flash ID + metadata recovery};
\node[box,below=of check] (ready) {Readiness gate: report READY or bounded diagnostic FAULT};
\node[box,below=of ready] (start) {START: allocate unique session, publish context, release acquisition};
\draw[->] (hal)--(rtos);\draw[->] (rtos)--(check);\draw[->] (check)--(ready);\draw[->] (ready)--(start);
\end{tikzpicture}
\caption{Proposed startup sequence and readiness gate.}\end{figure}
\begin{figure}[htbp]\centering\begin{tikzpicture}[>=Stealth,box/.style={draw=SEDSBlue,rounded corners,align=center,text width=28mm,minimum height=10mm,font=\small},lab/.style={font=\scriptsize,align=center,fill=white,inner sep=1pt}]
\node[box] (boot) at (0,0) {BOOT};
\node[box] (ready) at (5,0) {READY};
\node[box] (log) at (10,0) {LOGGING};
\node[box] (fault) at (0,-2.3) {FAULT};
\node[box] (maint) at (5,-2.3) {MAINTENANCE};
\node[box] (stop) at (10,-2.3) {STOPPING};
\node[box] (full) at (10,-4.3) {FULL};
\draw[->] (boot)--node[above,lab]{checks pass}(ready);
\draw[->] (ready)--node[above,lab]{START}(log);
\draw[->] (boot)--node[left,lab]{failure}(fault);
\draw[->] (log)--node[right,lab]{STOP}(stop);
\draw[->] (stop)--node[above,lab]{drained /\\verified}(maint);
\draw[->] (maint)--node[right,lab]{new session\\if healthy}(ready);
\draw[->] (log.east)--(12.2,0)--node[right,lab]{no erased\\space}(12.2,-4.3)--(full.east);
\draw[->] (full)--node[left,lab]{STOP}(stop);
\end{tikzpicture}

\small Any state may enter FAULT on an unrecoverable error. FULL preserves retained sessions; STOP still drains and closes.\normalsize
\caption{Operating states with storage and fault exits.}\end{figure}
BOOT performs HAL initialization and starts ThreadX. Platform object creation failures prevent logging. Data identifies flash and reconstructs session allocations. Sensor initializes devices; required-device failure prevents START. READY reports health and awaits explicit START. LOGGING creates a unique session and accepts both record types.

STOP establishes a producer cutoff. Sensor acknowledges its last sequence; processing drains all accepted earlier samples and acknowledges its last result. Only after both acknowledgments and the record queue drain does data verify pending writes and close the session. Bound the handshake to an initial 2 s budget; on timeout mark the session incomplete and report the last durable sequence rather than claiming success. Rework this budget if measured drain time requires it.

MAINTENANCE permits export and confirmed deletion of closed sessions. FULL disables additional persistence, counts rejected records and preserves retained data; acquisition and processing can continue for status. Reclaim space only through explicit maintenance deletion, then start a new session. FAULT reports its cause and prevents unsafe writes until a documented reset/recovery procedure succeeds. A USB disconnect releases leases and tokens; it does not automatically stop acquisition or delete a log.

\begin{figure}[htbp]\centering\begin{tikzpicture}[>=Stealth,font=\small]
\foreach \x/\name in {0/Sensor,4/Processing,8/Data,12/Flash}{
\node[draw=SEDSBlue,rounded corners,minimum width=23mm] at (\x,0) {\name};
\draw[dashed,gray] (\x,-0.3)--(\x,-5.8);}
\draw[->] (0,-0.8)--node[above,font=\scriptsize]{raw to process queue}(4,-0.8);
\draw[->] (0,-1.6)--node[above,font=\scriptsize]{independent raw to data queue}(8,-1.6);
\draw[->] (4,-2.4)--node[above,font=\scriptsize]{processed result}(8,-2.4);
\draw[->] (8,-3.2)--node[above,font=\scriptsize]{program body/CRC}(12,-3.2);
\draw[<-] (8,-4)--node[above,font=\scriptsize]{poll + verify}(12,-4);
\draw[->] (8,-4.8)--node[above,font=\scriptsize]{program commit}(12,-4.8);
\draw[<-] (8,-5.6)--node[above,font=\scriptsize]{verify then durable}(12,-5.6);
\end{tikzpicture}
\caption{Two independent raw deliveries and verified persistence.}\end{figure}
